\documentclass[12pt, a4paper]{article}
\usepackage[utf8]{inputenc}
\usepackage{mathptmx} % Times New Roman
\usepackage[margin=1in]{geometry}
\usepackage{setspace}

\title{\textbf{Integrating Artificial Intelligence into University Teaching Methods: A Study of Student Acceptance}}
\author{}
\date{}

\begin{document}

\maketitle
\vspace{-4em}

\begin{abstract}
\noindent Artificial intelligence continues to advance rapidly and influence many areas of society, including higher education. Traditional teaching methods still provide the core foundation for university instruction, but generative AI tools now offer new ways to support more interactive and engaging forms of learning.

Although these tools have become increasingly accessible, empirical research has yet to fully clarify the specific factors that shape students’ acceptance of AI or how such tools may alter established teaching practices. This study therefore examines the key determinants of artificial intelligence adoption among university students and considers its potential effects on conventional educational strategies.

The primary aim of the research is to evaluate the pedagogical value of incorporating artificial intelligence into university curricula and to construct a model capable of predicting students’ behavioral intentions toward using these tools in their academic work.

The study adopted a quantitative approach. Self-reported data were collected from a diverse sample of 500 undergraduate students through structured questionnaires. Structural equation modeling was then used to test hypothesized relationships involving technology acceptance, performance expectancy, and habit.

The findings from this research are expected to contribute useful insights for the integration of digital tools in higher education. A clearer understanding of the factors that drive student engagement with AI can help educators and administrators develop more effective approaches to implementing new teaching methods and creating a more dynamic learning environment.
\end{abstract}

\end{document}
